\documentclass[../Main.tex]{subfiles}

\begin{document}

\section{Recherche zum Leveldesign} \label{sec:Leveldesign}
Dieses Kapitel bündelt die recherchierten Designprinzipien thematisch in sechs Überkategorien. Sie stammen aus zwei Fachbüchern und vier GDC-Vorträgen, deren Auswahl wir in Abschnitt~\ref{sec:VorgehenLeveldesign} beschreiben. Wir beschränken uns dabei vor allem auf Prinzipien, die die Gegnerplatzierung und den Kampfraum betreffen, da diese im Zentrum unserer Erkenntnisfrage stehen. Diese werden besonders in Bezug auf das Hack-and-Slash-Genre verstanden. Wo möglich wird zusätzlich Bezug auf die Referenzspielanalyse (Kapitel~\ref{sec:Referenzspielanalyse}) genommen. Die unten genannten Designprinzipien variieren von abstrakten Denkmustern über praktisch anwendbare Methoden bis hin zu Gamedesign-Prozessen. Das liegt an der Vielfalt der Quellen und den verschiedenen Absichten der Autoren.
\begin{itemize}[leftmargin=1.5cm]
    \item \textcolor{red}{\textbf{Legende:} Roter Fliesstext braucht ihr für die Erkenntnisfrage nicht und kann raus. Rote Stichpunkte sind Notizen dazu.}
\end{itemize}

\subsection{Kernmechanik-zentriertes Design} \label{sec:Kernmechanik}

Ein Videospiel wird vor allem wegen seiner Kernmechanik gespielt. Diese ist das grundlegende Verb, das die Spielhandlung trägt.\autocite[110]{totten2014architectural} Bei \game{Super Mario Bros.} ist dies das Springen, bei \game{The Legend of Zelda} das Erkunden.\autocite{nintendo_supermariobros_1985,nintendo_zelda_1986} So liegt der Fokus der Spieler auf genau diesen Kernmechaniken, was nicht heisst, dass Spiele keine Zweit- oder Drittmechaniken haben können.

Das Leveldesign wird von den Mechaniken getrieben. Das Level ist demnach das Medium, durch das sich die Mechaniken entfalten.\autocite{gdc_tenprinciples2010} Deshalb ist es sinnvoll, jene Objektklasse, die zur Kernmechanik beiträgt, so breit wie möglich im Spiel zu verankern. Ist die Kernmechanik beispielsweise das Zerstören von Objekten, sollte möglichst vieles von dieser Klasse erben, sodass die Mechanik auf die eine oder andere Weise überall im Spiel auftaucht.\autocite{gdc_movingmountains2020}

Für \game{Project Polearm} heisst das, dass die Welt auf schnellen, aggressiven Combat ausgerichtet sein muss und dass beispielsweise alle Gegner von derselben grundlegenden Klasse erben.\footnote{Wie das bei uns umgesetzt ist, beschreibt der \textit{Project Polearm Technical Report}.} Da sich die Mechaniken im Level entfalten, bemessen wir auch unsere Räume daran, wie viel Bewegungsfreiraum sie dem Spieler für diesen Kampf lassen.

\subsection{Kampf-Encounter-Design} \label{sec:Encounter}
Encounter-Design hat viel mit der \textit{Lesbarkeit} eines Levels oder eines Spielabschnitts zu tun. \textcolor{red}{Unserer Erfahrung nach bewegen sich Spieler schneller und unbekümmerter durch offene und überschaubare Flächen, während Horrorspiele eher in engen und geschlossenen Räumlichkeiten stattfinden.} Ein Kampfraum sollte bereits vom Eingang aus lesbar sein: Ziel, Gegnerpositionen und Deckung müssen erkennbar sein, damit der Spieler eine Strategie fassen kann.\autocite{gdc_horizonquests2018} Je nachdem, wie übersichtlich ein Raum ist, muss ihn ein Spieler ein- oder mehrmals scannen. Bei Multiplayer-Spielen macht dies einen Skillunterschied aus, da geübte Spieler den Raum schneller und mit den richtigen Abständen scannen.\autocite[122--123]{totten2014architectural} Ein Beispiel dafür ist unserer Erfahrung nach \game{Counter-Strike},\autocite{valve_counterstrike_2000} in dem geübte Spieler einen Raum beim Betreten mit wenigen Blicken erfassen. Dass ein schnell erfassbarer Raum hilft, deckt sich mit unserer Beobachtung in \game{NieR: Automata}, wo die klar lesbare Anfangsaufstellung der Gegner sofort einen Überblick gibt (Frage~A3, Kapitel~\ref{sec:Referenzspielanalyse}). Umgekehrt lässt sich die Lesbarkeit unserer Ansicht nach auch bewusst einschränken: In \game{Ninja Gaiden Sigma Plus} tauchen neue Gegner häufig ausserhalb des Sichtfeldes auf, was die Spielhaftigkeit verbirgt (siehe Frage~A3, Kapitel~\ref{sec:Referenzspielanalyse}).

\textcolor{red}{Verschiedene Pfade ermöglichen es dem Spieler, seinen präferierten Spielstil auszuleben, da ein Pfad beispielsweise aggressives Verhalten und ein anderer das Schleichen vereinfacht.\autocite{gdc_tenprinciples2010}} Gute Kampfräume bieten dafür mehrere Routen, Flankiermöglichkeiten und Höhenunterschiede. Milde Barrieren wie eine Umkehr-Treppe zwingen den Spieler zudem zum kurzen Innehalten, bevor er in den Kampf stürmt.\autocite{gdc_horizonquests2018} \textcolor{red}{Des Weiteren kann man den Spieler unterschwellig leiten, indem man Pickups als Wegweiser streut. Dabei halten sich Spieler oft an helle Pfade und steigen gerne nach oben, während Senkungen oder Absperrungen den Spieler alarmieren.\autocite[240, 245, 251]{rogers2014levelup}}

Auch die Platzierung der Gegner beeinflusst das Laufverhalten des Spielers: Je nach Spiel meidet der Spieler Gegner oder greift sie bewusst an. So können Gegner selbst zur unsichtbaren Wand oder zum Wegweiser werden.\autocite[199]{totten2014architectural} Eine grosse Zahl oder besonders starke Gegner können den Spieler sogar aktiv in eine gewünschte Richtung drängen. Dieses Prinzip wird als \textit{Shrinking Fortress} bezeichnet und wirkt gerade in eigentlich grossen Räumen, weil Gegnerschwärme den Raum gefühlt verengen.\autocite[199--200]{totten2014architectural} Das deckt sich mit unserer Beobachtung zu \game{Ninja Gaiden Sigma Plus}, wo die Umzingelung durch mehrere Gegner ein einfaches Rückwärtsausweichen verhindert und den Spielfluss beschleunigt (Frage~D3, Kapitel~\ref{sec:Referenzspielanalyse}).

\textcolor{red}{Über die reine Platzierung hinaus prägt das koordinierte Verhalten der Gegner den Verlauf eines Encounters. In \game{Ninja Gaiden Sigma Plus} liessen sich drei Rollen beobachten: Nahkämpfer, die den Spieler in den Kampf ziehen, Flankierer, die ihn umkreisen, und Fernkämpfer, die ihn aus der Distanz unter Druck setzen (Frage~M2, Kapitel~\ref{sec:Referenzspielanalyse}). Erst dieses Zusammenspiel erzeugt unserer Beobachtung nach taktischen Druck.}

\begin{itemize}[leftmargin=1.5cm]
    \item \textcolor{red}{Rot markiert: Satz zu Horrorspielen -- hatte keine Quelle (jetzt als eigene Erfahrung formuliert) und wird für die Erkenntnisfrage nicht gebraucht.}
    \item \textcolor{red}{Rot markiert: Pfade fürs Schleichen und Wegführung über Pickups -- euer Kampfabschnitt ist linear und ohne Schleichen. Fällt der Satz zu den Pfaden weg, im folgenden Satz \enquote{dafür} streichen.}
    \item \textcolor{red}{\textbf{Letzter Absatz (Rollen in \game{Ninja Gaiden}) rot:} Das ist keine Leveldesign-Literatur, sondern eine Wiederholung aus Kapitel~\ref{sec:Referenzspielanalyse} (Frage~M2). Tabelle~\ref{tab:Designvorschlaege} verweist nicht mehr darauf, der Absatz kann also raus.}
\end{itemize}

\subsection{Räumliche Typologien} \label{sec:Raumtypologien}
Die räumliche Gestaltung lässt sich in \textit{Narrow Space}, \textit{Intimate Space} und \textit{Prospect Space} einteilen. Jeder dieser Typen symbolisiert für den Spieler eine andere Situation oder Emotion: Narrow Spaces stehen für Engpässe, Hinterhalte und Klaustrophobie. Intimate Spaces bilden den sicheren Standardfall, etwa einen Tutorialraum. Prospect Spaces sind offen und machen verwundbar, da Angriffe von überall möglich sind, weshalb sie oft als Bossräume dienen.\autocite[161--168]{totten2014architectural} Eine ähnliche Unterscheidung findet sich bei Rogers, der grosse Räume für grössere Gegnergruppen und enge Räume für Einzelkämpfe empfiehlt.\autocite[245]{rogers2014levelup} In \game{Ninja Gaiden Sigma Plus} konnten wir beobachten, dass die Räume im Verlauf gezielt schmaler wurden, was unserer Ansicht nach einem Übergang von Intimate zu Narrow Space entspricht (Frage~M3, Kapitel~\ref{sec:Referenzspielanalyse}).

\textcolor{red}{Das Entwerfen von Räumen lässt sich durch \textit{Molekül-Design} und \textit{Proximity-Diagramme} erleichtern. Dabei geht es nicht um genaue Abmessungen, sondern darum, welche Räume miteinander verbunden sind und wie der Spielfluss verlaufen soll.\autocite[168--169]{totten2014architectural} Zwei Beispiele solcher Diagramme zeigen Abbildung~\ref{fig:proximity} und Abbildung~\ref{fig:molekuel}.}

\begin{figure}[H]
    \centering
    \begin{minipage}{0.49\textwidth}
        \centering
        \includegraphics[height=5cm]{../Assets/Proximity-Diagram.png}
        \captionof{figure}{Proximity-Diagramm eines Multiplayer-Shooters, Quelle: Totten}
        \label{fig:proximity}
    \end{minipage}
    \hfill
    \begin{minipage}{0.49\textwidth}
        \centering
        \includegraphics[height=5cm]{../Assets/Molekühl-Design.png}
        \captionof{figure}{Molekül-Diagramm eines Spielabschnitts, Quelle: Totten}
        \label{fig:molekuel}
    \end{minipage}
\end{figure}
\textcolor{red}{Kevin Lynch unterteilt städtische Räume in fünf Elemente: Landmarks, Paths, Nodes, Edges und Districts. Dieses Modell war ursprünglich für grosse, offene Welten gedacht. Für einen linearen Kampfabschnitt sind davon vor allem \textit{Edges} und \textit{Nodes} nützlich. Edges sind Grenzen oder Übergänge, die einen Ortswechsel signalisieren, etwa durch einen Wechsel von Textur, Bewuchs oder Baustil. Nodes sind Entscheidungspunkte, an denen sich Pfade kreuzen. Beide helfen dabei, die Übergänge zwischen den einzelnen Abschnitten eines Kampfareals klar zu markieren, sodass der Spieler erkennt, wann er einen neuen Abschnitt betritt.\autocite[182--184]{totten2014architectural}}

\begin{itemize}[leftmargin=1.5cm]
    \item \textcolor{red}{Rot markiert: Molekül- und Proximity-Diagramme sowie Lynch betreffen die Planung ganzer Level, nicht Gegner. Nur behalten, wenn ihr diese Diagramme selbst verwendet habt. Werden sie gestrichen, auch die beiden Abbildungen entfernen.}
    \item \textcolor{red}{Falls ihr die Abbildungen behaltet: In der Bildunterschrift reicht \enquote{Quelle: Totten} nicht, es braucht einen Quellenbeleg mit Seitenzahl (Zitierregeln). Lynch ist ein Zitat aus zweiter Hand: \enquote{Totten stützt sich dabei auf Kevin Lynch}.}
\end{itemize}

\subsection{Spannungskurve, Pacing \& Überraschung} \label{sec:Pacing}
Spannung entsteht nicht nur durch Erzählung oder Musik, sondern bereits durch das Levellayout selbst. Spiele setzen dabei oft abwechselnd auf Phasen hoher Intensität und tiefer Entspannung. Dieses Muster wird als \textit{Rollercoaster-Pacing} bezeichnet und wiederholt sich fraktalähnlich auf der Ebene einzelner Abschnitte wie auch des gesamten Spiels.\autocite{gdc_tenprinciples2010} Es deckt sich mit der Wellenstruktur, die wir in \game{Ninja Gaiden Sigma Plus} beobachtet haben, wo die Spannung mit jeder Welle steigt und zwischen den Wellen kurz abfällt (Frage~D1, Kapitel~\ref{sec:Referenzspielanalyse}).

\textcolor{red}{Besonders das Unvorhergesehene reizt die Spieler, denn ein zu regelmässiges Pacing wird schnell langweilig. Der sogenannte \textit{Designer's Curse} beschreibt dabei ein Problem der Entwickler selbst: Sie durchschauen ihre eigenen Muster so schnell, dass sie deren Überraschungswirkung kaum noch einschätzen können. Das Gegenmittel ist, etablierte Muster gezielt zu brechen, etwa wie in \game{Dead Space 2},\autocite{visceral_deadspace2_2011} wo ein Abschnitt bewusst lange ohne Gegnerkontakt bleibt, bevor die Bedrohung eskaliert.\autocite{gdc_tenprinciples2010}}

\textcolor{red}{Ein flexibles Pacing lässt sich über mehrere Pfade erreichen. Wichtig ist dabei, dass Risiko und Belohnung jedes Pfades klar kommuniziert werden, sodass der Spieler seine Schwierigkeit selbst wählt, statt sie blind über ein Menü festzulegen, etwa über sichtbar markierte, aber riskante Abkürzungen wie in \game{Burnout}.\autocite{criterion_burnout_2001,gdc_tenprinciples2010}}

\textcolor{red}{\textit{Backtracking} ist ein umstrittenes Thema. Für die einen hat es einen spielerischen Wert, für die anderen streckt es das Spiel unnötig. Ist Backtracking Teil des Spiels, sollte sich der Rückweg spielerisch vom Hinweg unterscheiden.\autocite{gdc_tenprinciples2010}}

Übergeordnet sollte beim Gestalten eines Levels immer zuerst vom gewünschten Gefühl ausgegangen werden, danach werden die Mechaniken und Hindernisse ausgewählt und eingebaut, die dieses Gefühl erzeugen.\autocite{gdc_tenprinciples2010}

\begin{itemize}[leftmargin=1.5cm]
    \item \textcolor{red}{Rot markiert: Pacing über mehrere Pfade (\game{Burnout}) -- euer Kampfabschnitt hat nur einen Weg.}
    \item \textcolor{red}{Rot markiert: \textit{Designer's Curse} und \game{Dead Space 2} -- wird in keinem Designvorschlag verwendet.}
    \item \textcolor{red}{Rot markiert: Backtracking -- nur behalten, wenn ihr damit begründen wollt, dass Raum~12 derselbe Raum wie Raum~3 ist, aber anders bespielt wird (Abschnitt~\ref{sec:Kampfabschnitt}). Sonst raus.}
\end{itemize}

\subsection{Spielergeführtes Lernen / Scaffolding} \label{sec:Scaffolding}
Das \textit{Scaffolding} beschreibt als Leveldesignprinzip, wie neue Mechaniken in einem Level am besten eingeführt werden.\autocite[114]{totten2014architectural} Grundsätzlich sollte pro Level nur eine einzige Mechanik neu hinzukommen.\autocite{gdc_tenprinciples2010} Sie wird zunächst eigenständig und gefahrlos gezeigt, idealerweise ganz ohne erklärenden Text, und erst danach mit weiteren Mechaniken verflochten, was als \textit{Braiding} bezeichnet wird.\autocite{gdc_movingmountains2020} \textcolor{red}{Konkrete Muster dafür sind das Zelda-Dungeon-Modell, bei dem jeder Raum eine neue Anwendung eines Items lehrt und der Bosskampf die Beherrschung prüft, sowie die Bethesda-Schleife aus Lernen, Spielen, Herausforderung und Überraschung.\autocite{gdc_tenprinciples2010}}

So soll beispielsweise in \game{Project Polearm} ein Fernkämpfer isoliert eingeführt werden, bevor er mit einem Nahkämpfer kombiniert auftritt. Damit lässt sich der taktische Druck, den wir in \game{Ninja Gaiden Sigma Plus} bei der Kombination verschiedener Gegnertypen beobachtet haben, schrittweise aufbauen (Frage~D2, Kapitel~\ref{sec:Referenzspielanalyse}).

\begin{itemize}[leftmargin=1.5cm]
    \item \textcolor{red}{Rot markiert: Zelda-Dungeon-Modell und Bethesda-Schleife -- für eure gestaffelte Einführung der Gegner nicht nötig.}
\end{itemize}

\subsection{Prototyping \& Iteration} \label{sec:Prototyping}
Level sollten nicht in einem Durchgang entworfen, sondern schrittweise entwickelt werden. Grundlegende Fragen lassen sich bereits vor jeder 3D-Geometrie klären, etwa mit nicht-digitalen Prototypen auf Papier oder mit Spielfiguren. Danach folgt üblicherweise die \textit{Greybox}, eine einfache Blockgeometrie, mit der Raumgrössen und Distanzen getestet werden, bevor die grafische Ausarbeitung beginnt.\autocite[124--126]{totten2014architectural}

\textcolor{red}{Ein konkreter Ablauf dafür führt vom Konzept über einen \textit{Walkthrough}, eine schriftliche Beschreibung des Levelverlaufs samt Draufsicht-Karte, zur \textit{Blockworld}. Dabei weicht die Blockworld oft stärker vom ursprünglichen Plan ab als erwartet, weil sich Sichtlinien, Kameraverdeckung und tatsächliche Raumgrössen erst im 3D-Raum beurteilen lassen. Als Planungswerkzeug davor eignen sich \textit{Bubble-Diagramme}, die Räume als Blasen und deren Beziehungen statt genauer Geometrie festhalten, ähnlich den Proximity-Diagrammen aus Abschnitt~\ref{sec:Raumtypologien}. Beim Bauen der Blockworld empfiehlt es sich zudem, von einem einzelnen markanten Kernpunkt aus zu wachsen, statt das ganze Level gleichzeitig anzulegen.\autocite{gdc_leveldesignaday2020}}

Wie zentral diese iterative Arbeitsweise ist, zeigt Valve am Beispiel von \game{Half-Life 2}\autocite{valve_halflife2_2004}: Die Level durchliefen mehrere vollständige Überarbeitungsdurchgänge, und die Entwickler bezeichnen ebendiese mehrfachen Iterationen rückblickend als einen Schlüssel zum Erfolg des Spiels, da sie deutlich bessere Designentscheidungen ermöglichten.\autocite{jacobson2005postmortem} Für \game{Project Polearm} entspricht dieser Ablauf unserem eigenen Vorgehen, den Kampfabschnitt zuerst als Greybox zu bauen und iterativ zu verfeinern, bevor Grafik und Politur folgen.

\begin{itemize}[leftmargin=1.5cm]
    \item \textcolor{red}{Rot markiert: Walkthrough, Blockworld und Bubble-Diagramme -- zu detailliert. Für den Verweis aus Abschnitt~\ref{sec:VorgehenUmsetzung} reichen der erste Absatz (Greybox) und \game{Half-Life 2}.}
    \item \textcolor{red}{Hinweis: 4.6 betrifft eure Arbeitsweise, nicht die Erkenntnisfrage. Wenn ihr kürzen müsst, könnte der Abschnitt auch ganz raus; dann den Verweis in Abschnitt~\ref{sec:VorgehenUmsetzung} entfernen.}
\end{itemize}

\subsection{Erkenntnisse für Project Polearm} \label{sec:ZwischenfazitLD}

Im Folgenden fassen wir kurz zusammen, welche Erkenntnisse aus diesem Kapitel für unsere Erkenntnisfrage und unseren Prototyp wichtig sind. Sie sind vor allem für den Teil der Erkenntnisfrage relevant, in dem es um die Gegnerplatzierung und das Leveldesign geht, während die Analyse der Referenzspiele (Kapitel~\ref{sec:Referenzspielanalyse}) vor allem das Gegnerverhalten betrifft.

Als übergeordneten Grundsatz gehen wir bei der Gestaltung unseres Prototyps vom gewünschten Gefühl aus und gestalten den Kampf danach (Abschnitt~\ref{sec:Pacing}). Für unser Spiel stellen wir uns einen schnellen und aggressiven Kampf vor, in dem sich der Spieler grundsätzlich überlegen fühlt (Abschnitte~\ref{sec:Genredefinition} und~\ref{sec:Vergleich}). Daraus folgt für uns, dass wir Räume, Hindernisse und Gegner so entwerfen müssen, dass sie dieses Gefühl unterstützen oder hervorrufen. Die konkreten Werkzeuge, mit denen wir diese Gefühle beeinflussen wollen, sind die folgenden.

Erstens muss ein Raum schnell lesbar sein, damit sich der Spieler rasch eine Strategie zurechtlegen kann (Abschnitt~\ref{sec:Encounter}).

Zweitens sind die Raumtypen Narrow, Intimate und Prospect Space besonders wichtig (Abschnitt~\ref{sec:Raumtypologien}). Wenn wir sie gezielt mischen, können wir sowohl das Pacing des Spiels als auch die Gefühle der Spieler lenken. Dafür müssen wir uns aber bewusst sein, wie die einzelnen Raumtypen wirken. So muss etwa die Anzahl der Gegner an den Raumtyp angepasst werden, wie es auch Rogers empfiehlt.\autocite[245]{rogers2014levelup}

Nicht nur die Grösse der Räume, sondern auch die Gegner selbst können den Spieler dazu zwingen, sich anders zu bewegen. Gegnergruppen können einen Raum sogar gefühlt verengen (\textit{Shrinking Fortress}, Abschnitt~\ref{sec:Encounter}). Damit haben wir unserer Ansicht nach ein weiteres Mittel, um die Bewegungen der Spieler vorherzusehen und abzuschätzen, welche Gefühle ein Raum bei ihnen auslöst (Abschnitt~\ref{sec:Raumtypologien}).

Drittens nutzen wir das Rollercoaster-Pacing, das Anspannung und Erholung abwechselt und sich fraktalähnlich auf mehreren Ebenen gleichzeitig wiederholt (Abschnitt~\ref{sec:Pacing}). Damit ist gemeint, dass sowohl der gesamte Spannungsbogen eines Spiels als auch die einzelnen Kampfabschnitte einem solchen Verlauf folgen.

Viertens sind das Scaffolding- und das Braiding-Prinzip für unseren Prototyp zentral (Abschnitt~\ref{sec:Scaffolding}). Übertragen auf Gegner bedeutet das, eine neue Gegnervariante zuerst einzeln und gefahrlos vorzustellen und sie erst danach nach und nach in Gruppen und Kombinationen ins Spiel einzubinden. Auch den Grundsatz, dass pro Level nur eine neue Mechanik hinzukommen soll, möchten wir auf unseren Prototyp übertragen, und zwar auf einzelne Spielabschnitte statt auf ganze Level.

All diese Erkenntnisse fliessen zusammen mit jenen aus der Analyse der Referenzspiele (Abschnitt~\ref{sec:Vergleich}) in unseren Prototyp und in die Beantwortung unserer Erkenntnisfrage ein. Im nächsten Kapitel zeigen wir, wie wir dabei konkret vorgehen.


\end{document}
